\documentclass[10pt,a4paper]{amsart}
\usepackage[margin=19mm]{geometry}
\usepackage{amsmath,amssymb,mathtools}
\usepackage[scr=rsfs,cal=euler]{mathalfa}
\usepackage{xfrac}
\usepackage[unicode,hidelinks,bookmarks=false]{hyperref}
\newcommand{\ie}{i.e.\ }
\newcommand{\cf}{cf.\ }
\newcommand{\resp}{respectively\ }
\newcommand{\Subsection}{\subsection}
\providecommand{\nobreakdash}{\nobreak\hbox{-}}
\setlength{\emergencystretch}{2em}
\multlinegap0pt
\usepackage{tikz}
\usetikzlibrary{cd}
\usepackage{amssymb}
\usepackage{xcolor}
\usepackage{enumerate}
\usepackage{mathtools}
\usepackage{bm}
\usepackage{cancel}
\usepackage[normalem]{ulem}
\usepackage[shortlabels]{enumitem}
\usepackage{tikz}
\newtheorem{lemma}{Lemma}
\title{A necessary condition for liftings of positive characteristic varieties with finite fundamental groups}
\author{Ruida Di, Runjie Hu, Siqing Zhang}
\date{Source: arXiv:2506.07914v3 (CC BY 4.0)}
\begin{document}
\maketitle

\noindent\emph{Source and licence.} A necessary condition for liftings of positive characteristic varieties with finite fundamental groups. Version arXiv:2506.07914v3 (CC BY 4.0), distributed by arXiv under the Creative Commons Attribution 4.0 International licence, http://creativecommons.org/licenses/by/4.0/, as recorded in the arXiv licence metadata for this exact version and shown on the abstract page https://arxiv.org/abs/2506.07914v3. Full licence notice: \texttt{licenses/arxiv-2506.07914-CC-BY-4.0.txt}.\par\smallskip

\noindent\textit{Editorial scope.} Counted unit: the article's lemma lemma (source0.tex lines 530--532). The article's complete original proof of it is delivered with it (source0.tex lines 534--544, one \texttt{proof} environment of 11 lines). No mathematics is written, repaired or re-derived: the statement and the proof are the article's own lines, in the article's own notation, and no retained line is substituted or re-flowed.  Retained uncounted as the source-local context the proof needs: lines 491--493.  Nothing was omitted from any retained span, so the package carries no omission record.  No retained line contains a citation, so the wrapper carries no bibliography and no bibliography entry.  The retained lines contain the article's own diagram code, which the wrapper typesets with the standard tikz packages; no image file is delivered and the package declares no asset dependency.  Editorial formatting only, in addition: an AMS \texttt{amsart} preamble that restates the article's own declarations and macros used by the retained lines, namely the environments \texttt{lemma} on separate counters, together with the standard packages those lines need.  The retained environments print in this document as Lemma 1 in order of appearance, whereas the article prints the same items with the numbering of its own full document.  The complete native source of the article as distributed in the arXiv e-print for this exact version is bundled unchanged as \texttt{sources/179567/source0.tex}.
\begin{lemma}\label{Lem: Commutativity of universal cover and l-completion}
Let $Y$ be a connected CW complex with $\pi_1(Y)$ finite. Then the canonical map $(\widetilde{Y})_l\rightarrow (Y_l)^{\sim}$ is a homotopy equivalence.
\end{lemma}

\begin{lemma}
Let $Z$ be a connected CW complex with $\pi=\pi_1(Z)$ a finite $l$-group. Assume that either $Z$ is finite type or $l$-local finite type. Then there is a natural homotopy equivalence $Z_l\rightarrow Z'$, where $Z'$ is the fiberwise $l$-completion of $Z$.
\end{lemma}

\begin{proof}
Consider the following diagram the fiberwise $l$-completion of $Z$.
\[
\begin{tikzcd}
   \widetilde{Z} \arrow[r] \arrow[d] &  Z \arrow[r] \arrow[d] & K(\pi,1) \arrow[d,phantom, sloped, "="] \\
   \widetilde{Z'}\simeq(\widetilde{Z})_l \arrow[r] &  Z' \arrow[r] & K(\pi,1)
\end{tikzcd}
\]
By the finite-type assumption, $Z'$ is $l$-complete finite type and $\pi_1(Z')$ is a finite $l$-group. Therefore, $\widetilde{Z'}$ is $l$-profinite complete. So $\pi_q(Z')\cong \pi_q(\widetilde{Z'})\cong \pi_q(((Z')_l)^{\sim})\cong \pi_q((Z')_l)$ for any $q\geq 2$ due to Lemma \ref{Lem: Commutativity of universal cover and l-completion}.
Hence $Z'$ is also $l$-profinite complete. Then the map $Z\rightarrow Z'$ factors through $Z_l$. This induces a map $(Z_l)^{\sim}\rightarrow \widetilde{Z'}$ and this map is a homotopy equivalence due to Lemma \ref{Lem: Commutativity of universal cover and l-completion}. 
\end{proof}

\end{document}
